\documentclass[10pt, letterpaper]{article}

% Packages:
\usepackage[
    ignoreheadfoot, % set margins without considering header and footer
    top=2 cm, % seperation between body and page edge from the top
    bottom=2 cm, % seperation between body and page edge from the bottom
    left=2 cm, % seperation between body and page edge from the left
    right=2 cm, % seperation between body and page edge from the right
    footskip=1.0 cm, % seperation between body and footer
    % showframe % for debugging 
]{geometry} % for adjusting page geometry
\usepackage{titlesec} % for customizing section titles
\usepackage{tabularx} % for making tables with fixed width columns
\usepackage{array} % tabularx requires this
\usepackage[dvipsnames]{xcolor} % for coloring text
\definecolor{primaryColor}{RGB}{0, 79, 144} % define primary color
\usepackage{enumitem} % for customizing lists
\usepackage{fontawesome5} % for using icons
\usepackage{amsmath} % for math
\usepackage[
    pdftitle={Isaac's CV},
    pdfauthor={Isaac Vazquez},
    pdfcreator={LaTeX with RenderCV},
    colorlinks=true,
    urlcolor=primaryColor
]{hyperref} % for links, metadata and bookmarks
\usepackage[pscoord]{eso-pic} % for floating text on the page
\usepackage{calc} % for calculating lengths
\usepackage{bookmark} % for bookmarks
\usepackage{lastpage} % for getting the total number of pages
\usepackage{changepage} % for one column entries (adjustwidth environment)
\usepackage{paracol} % for two and three column entries
\usepackage{ifthen} % for conditional statements
\usepackage{needspace} % for avoiding page brake right after the section title
\usepackage{iftex} % check if engine is pdflatex, xetex or luatex

% Ensure that generate pdf is machine readable/ATS parsable:
\ifPDFTeX
    \input{glyphtounicode}
    \pdfgentounicode=1
    % \usepackage[T1]{fontenc} % this breaks sb2nov
    \usepackage[utf8]{inputenc}
    \usepackage{lmodern}
\fi



% Some settings:
\AtBeginEnvironment{adjustwidth}{\partopsep0pt} % remove space before adjustwidth environment
\pagestyle{empty} % no header or footer
\setcounter{secnumdepth}{0} % no section numbering
\setlength{\parindent}{0pt} % no indentation
\setlength{\topskip}{0pt} % no top skip
\setlength{\columnsep}{0cm} % set column seperation
\makeatletter
\let\ps@customFooterStyle\ps@plain % Copy the plain style to customFooterStyle

\titleformat{\section}{\needspace{4\baselineskip}\bfseries\large}{}{0pt}{}[\vspace{1pt}\titlerule]

\titlespacing{\section}{
    % left space:
    -1pt
}{
    % top space:
    0.05 cm
}{
    % bottom space:
    0.2 cm
} % section title spacing

\renewcommand\labelitemi{$\circ$} % custom bullet points
\newenvironment{highlights}{
    \begin{itemize}[
        topsep=0.10 cm,
        parsep=0.10 cm,
        partopsep=0pt,
        itemsep=0pt,
        leftmargin=0.4 cm + 10pt
    ]
}{
    \end{itemize}
} % new environment for highlights

\newenvironment{highlightsforbulletentries}{
    \begin{itemize}[
        topsep=0.10 cm,
        parsep=0.10 cm,
        partopsep=0pt,
        itemsep=0pt,
        leftmargin=10pt
    ]
}{
    \end{itemize}
} % new environment for highlights for bullet entries


\newenvironment{onecolentry}{
    \begin{adjustwidth}{
        0.2 cm + 0.00001 cm
    }{
        0.2 cm + 0.00001 cm
    }
}{
    \end{adjustwidth}
} % new environment for one column entries

\newenvironment{twocolentry}[2][]{
    \onecolentry
    \def\secondColumn{#2}
    \setcolumnwidth{\fill, 4.5 cm}
    \begin{paracol}{2}
}{
    \switchcolumn \raggedleft \secondColumn
    \end{paracol}
    \endonecolentry
} % new environment for two column entries

\newenvironment{header}{
    \setlength{\topsep}{0pt}\par\kern\topsep\centering\linespread{1.5}\selectfont
}{
    \par\kern\topsep
} % new environment for the header

\newcommand{\placelastupdatedtext}{% \placetextbox{<horizontal pos>}{<vertical pos>}{<stuff>}
  \AddToShipoutPictureFG*{% Add <stuff> to current page foreground
    \put(
        \LenToUnit{\paperwidth-2 cm-0.2 cm+0.05cm},
        \LenToUnit{\paperheight-1.0 cm}
    ){\vtop{{\null}\makebox[0pt][c]{
        \small\color{gray}\textit{Last updated in September 2024}\hspace{\widthof{Last updated in September 2024}}
    }}}%
  }%
}%

% save the original href command in a new command:
\let\hrefWithoutArrow\href

% new command for external links:
\renewcommand{\href}[2]{\hrefWithoutArrow{#1}{\ifthenelse{\equal{#2}{}}{ }{#2 }\raisebox{.15ex}{\footnotesize \faExternalLink*}}}


\begin{document}
    \linespread{0.833}\selectfont % 10pt text on 10pt leading = 1.0 line spacing
    \newcommand{\AND}{\unskip
        \cleaders\copy\ANDbox\hskip\wd\ANDbox
        \ignorespaces
    }
    \newsavebox\ANDbox
    \sbox\ANDbox{\hspace{1em}{\small{\textbar}}\hspace{1em}}

    \begin{header}
        \textbf{\fontsize{17pt}{17pt}\selectfont Isaac Alexander Vazquez}

        \vspace{0.05cm}
        \mbox{\href{tel:+18053300059}{\color{black}{\footnotesize\faPhone*}\hspace*{0.13cm}(805) 330-0059}}%
        \kern 0.25cm%
        \AND%
        \normalsize
        \mbox{\href{mailto:ialexandervazquez5@gmail.com}{\color{black}{\footnotesize\faEnvelope[regular]}\hspace*{0.13cm}ialexandervazquez5@gmail.com}}%
        \kern 0.25cm%
        \AND%
        \kern 0.25cm%
        \mbox{\href{https://www.linkedin.com/in/ivazquezalexander/}{\color{black}{\footnotesize\faLinkedinIn}\hspace*{0.13cm}linkedin.com/in/ivazquezalexander/}}%
        \kern 0.25cm%
        \AND%
        \kern 0.25cm%
        \mbox{\href{https://github.com/ivazquez805}{\color{black}{\footnotesize\faGithub}\hspace*{0.13cm}ivazquez805}}%

    \end{header}

    \vspace{0.05 cm - 0.05 cm}



\section{Education}

\begin{twocolentry}{
\textit{San Luis Obispo, CA} \\
\textit{Expected Dec 2026}}
\textbf{California Polytechnic State University} \\
\textit{Bachelor of Science in Computer Engineering}
\end{twocolentry}

\begin{onecolentry}
\begin{highlights}
    \item \textbf{Relevant Courses:} Digital and Computer Design, Microcontrollers and Embedded Applications, Real-time Operating Systems, High-Performance Embedded Systems, Autonomous Mobile Robotics, Hardware Security
\end{highlights}
\end{onecolentry}


\section{Skills}

\begin{onecolentry}
\begin{highlights}
    \item \textbf{Programming languages:} C, C++, C\#, Python, Java, SystemVerilog, RISC-V Assembly
    \item \textbf{Hardware Platforms:} FreeRTOS, Linux, ROS2 Humble, Oscilloscope, Multimeter, GazeboSim, Foxglove
    \item \textbf{Microcontrollers:} STM32 ARM-Cortex Microcontrollers, NVIDIA Jetson Nano Dashboard, ESP32
    \item \textbf{Languages:} Spanish (fluent)
\end{highlights}
\end{onecolentry}
   
\section{Technical Projects}

\textbf{Amiga Robot Driver Configurations} \textbar{} \textit{ROS2, NVIDIA Jetson, CAN Bus} \hfill \textit{Aug 2026}
\vspace{0.08 cm}
\begin{onecolentry}
\begin{highlights}
    \item Joined an existing four-wheeled Agriculture robotics project and re-evaluated drivers built by a previous team to assess performance.
    \item Sniffed CAN packets between a PC and a Jetson Nano, analyzing frames with ROS2 Linux commands and oscilloscope measurements of voltage changes.
    \item Documented debugging processes for future collaborators and created launch files to start multiple ROS2 nodes concurrently.
\end{highlights}
\end{onecolentry}
\vspace{0.14 cm}

\textbf{RTOS Traffic Light System} \textbar{} \textit{FreeRTOS, Embedded Systems, STM32} \hfill \textit{June 2026}\vspace{0.08 cm}
\begin{onecolentry}
\begin{highlights}
    \item Re-architected a bare-metal traffic light system into a multi-task FreeRTOS application handling button input and asynchronous interrupts.
    \item Replaced polling and blocking delays with interrupt-driven button handling, using FreeRTOS APIs for precise timing and inter-task resource management.
    \item Debugged over UART using PuTTY and used software timers with callbacks to reduce task overhead.
\end{highlights}
\end{onecolentry}
\vspace{0.14 cm}

\textbf{Function Generator} \textbar{} \textit{Embedded Systems, Timer Interrupts, SPI, STM32} \hfill \textit{Oct 2025}
\vspace{0.08 cm}
\begin{onecolentry}
\begin{highlights}
    \item Designed and debugged an STM32L476RG system that generates AC waveforms with user-selectable shape, duty cycle, and frequency.
    \item Implemented SPI communication (master mode) to drive an external 12-bit DAC for waveform output.
    \item Achieved 100-500 Hz output from keypad input using polling and timer interrupts, verified on an oscilloscope.
\end{highlights}
\end{onecolentry}
\vspace{0.14 cm}

\textbf{RISC-V OTTER Microcontroller with Interrupts} \textbar{} \textit{FPGA, Vivado, RISC-V Assembly} \hfill \textit{May 2024}\vspace{0.08 cm}
\begin{onecolentry}
\begin{highlights}
    \item Designed and debugged a 32-bit RISC-V (RV32I) processor in SystemVerilog using Vivado, including an ALU module, Program Counter and Control Decoder Finite State machine running RISC-V Assembly programs
    \item Implemented interrupt support using CSR registers and an MRET instruction, replacing polling so the CPU can service button and switch events without wasting clock cycles
    \item Verified with waveform-simulation in Vivado, then tested on a Xilinx Basys3 board with custom Assembly programs, such as a switch-activated LED counter
\end{highlights}
\end{onecolentry}
\vspace{0.14 cm}

\section{Work Experience}
\vspace{0.14 cm}
% Cal Poly ITS
\begin{twocolentry}{
\textit{San Luis Obispo, CA} \\   
\textit{May 2025 – Present}}
\textbf{ITS Student Assistant}

\textit{Cal Poly Information Technology Services}
\end{twocolentry}

\vspace{0.10 cm}
\begin{onecolentry}
\begin{highlights}
    \item Provide technical support to students, faculty, and staff across hardware, software, and network issues, resolving about 10 issues per week.
    \item Image and deploy 5+ Windows and macOS devices for faculty and staff, installing applications and configuring campus network connectivity IPV4.
    \item Document troubleshooting steps from client tickets in Jira, creating a reference used by student assistants and new hires.
\end{highlights}
\end{onecolentry}

\vspace{0.2 cm}

\section{Involvements}
\vspace{0.14 cm}
% SHPE
\begin{twocolentry}{
\textit{San Luis Obispo, CA} \\
\textit{Aug 2022 -- Present}}
\textbf{Society of Hispanic Professional Engineers (SHPE)} \\
\textit{Member}
\end{twocolentry}
\vspace{0.2 cm}

\end{document}